\documentclass[11pt,a4paper]{article}
\usepackage[utf8]{inputenc}
\usepackage[T1]{fontenc}
\usepackage{lmodern}
\usepackage[margin=1in]{geometry}
\usepackage{hyperref}
\usepackage{enumitem}
\usepackage{xcolor}
\hypersetup{colorlinks=true, linkcolor=blue!60!black, urlcolor=blue!60!black}
\setlength{\parskip}{0.5em}
\setlength{\parindent}{0pt}

\title{\Large\bfseries Fate-Engineered Hair Dye --- Explained Simply}
\author{Based on: Green\_Sustainable\_Permanent\_Coloration\_Full\_Paper.tex}
\date{\today}

\begin{document}
\maketitle

\begin{quote}
\emph{Normal permanent hair dye lasts well, but it can cause allergies and it pollutes water. This paper proposes a dye that is \textbf{stable in your hair, self-destructs on your skin, and breaks down in wastewater}.}
\end{quote}

\section{Why Permanent Dye Is Hard to Replace}
Over 50\% of women and a growing share of men use oxidative hair color. Permanent dye dominates because only it covers grey durably in a single step.

It works in 3 steps:
\begin{enumerate}[leftmargin=1.5em]
\item \textbf{Swell and open:} The cream (pH 9--11) lifts the hair cuticle and swells the hair so dye can enter.
\item \textbf{Penetrate small, trap large:} Tiny color building-blocks (100--150 Da) go \emph{inside} the hair. Peroxide then glues them into big colored molecules ($>$500 Da) too big to wash out. That trap is what makes it permanent.
\item \textbf{Couple to color:} Primary intermediate + coupler $\rightarrow$ colored dye.
\end{enumerate}

Nothing else gives that same shade range + wash and light fastness.

The problem: those same building-blocks cause trouble elsewhere:
\begin{itemize}[leftmargin=1.5em]
\item \textbf{On skin:} leftover dye can stick to skin proteins and trigger allergy / sensitization.
\item \textbf{In rivers:} the big colored molecules that wash out don't break down, so they persist as pollution.
\end{itemize}

We want a dye that lasts AND disappears. Sounds impossible.

\section{The New Idea: Fate Engineering}
Instead of finding one perfect chemical, design a chemical that \textbf{behaves differently in 3 places}:

\subsection*{In Hair = STABLE}
Stays locked in, keeps color for 28 shampoos. Target: color change $\Delta E \leq 3$ after 28 shampoos, $\geq$95\% grey coverage.

\subsection*{On Skin = CLEAVABLE}
Your skin has enzymes called esterases. The new dye has a weak link --- an ester or carbonate side-chain --- that those enzymes cut.

Once cut, it becomes water-loving (polar) and washes off instead of sticking to proteins and causing allergy.

Think of it like: a key that melts as soon as it touches skin, so it can't unlock an allergic reaction. Target: half-life $\leq$ 60 minutes on skin-like models.

\subsection*{In Wastewater = DEGRADABLE}
The final color has another weak link, a ``trigger link'' placed away from the color part so it doesn't change the color.

In wastewater with bacteria / different pH, that link breaks and the color falls apart so bacteria can eat it. Target: $\geq$60\% breakdown in 28 days in OECD 301B eco-test.

\section{How They Plan to Test It}
\textbf{Important: This is only a proposal. No lab results yet.}

Plan is 12 months. Make 5--8 versions of the new dye and compare them to the current best low-allergy dye, called ME-PPD.

\textbf{Phase 1 (Months 1--4): Chemistry + Safety in Lab}
\begin{itemize}[leftmargin=1.5em]
\item Make the new molecules
\item Does it still make color fast enough? Target: within 25\% of ME-PPD speed
\item Does it survive pH 9--10 in the cream, but break on skin enzymes?
\item Allergy lab tests: DPRA, KeratinoSens, h-CLAT (no animals)
\end{itemize}

\textbf{Phase 2 (Months 5--9): Real Hair Tests}
\begin{itemize}[leftmargin=1.5em]
\item Test on natural, bleached, and 100\% white hair
\item Color yield, 28-shampoo fastness, light fastness
\item Hair strength and combing force
\end{itemize}

\textbf{Phase 3 (Months 10--12): Eco Fate}
\begin{itemize}[leftmargin=1.5em]
\item OECD 301B biodegradation test on precursors AND final dyes separately
\item Identify breakdown products
\item Screening aquatic toxicity
\end{itemize}

If a candidate fails any gate, it gets dropped.

\section{Why It's New}
Old fixes only solved one problem at a time:
\begin{itemize}[leftmargin=1.5em]
\item \textbf{ME-PPD:} less allergy, but still causes reaction in $\sim$30\% of allergic people, and still pollutes. Not designed to degrade.
\item \textbf{Amino-acid cream instead of ammonia:} less smell, but doesn't fix dye pollution.
\item \textbf{Melanin / natural coating:} safer, no aromatic amines, but very few colors and mostly surface coating that washes off.
\item \textbf{Resorcinol-free:} fixes coupler issue, not precursor persistence.
\end{itemize}

Nobody has tried to design the dye itself to have 3 different intentional fates on purpose. That's the new idea here.

The paper does NOT claim a new alkali, new oxidant, or full shade chart. Arginine alkali and normal peroxide are just used as controls. Focus is only on precursor + chromophore design.

\section{Risks}
\begin{enumerate}[leftmargin=1.5em]
\item The weak ester link might break too early in the alkaline cream before it even reaches hair. Must test this first.
\item Changing the molecule might slow down color formation.
\item Lab allergy tests don't replace human clinical tests.
\item The eco-trigger link might shift the color.
\item New dyes need full regulatory approval --- out of scope for 12 months.
\end{enumerate}

\section{If Your Friend Asks ``So Did It Work?''}
Answer: We don't know yet. This paper is the blueprint for \emph{how} to try it, with success criteria defined in advance.

It asks: Can we get stable-in-hair + safe-on-skin + degradable-in-water without losing performance vs ME-PPD? That test hasn't been run yet.

\end{document}
